\SourceLeaf{375}%
lung der gültige wäre, \(DC=CL\) sein, in Widerspruch mit dem vorher gefundenen. Wir haben, wie wir glauben, bei diesem Moment des Beweises, das worauf es ankommt, noch etwas klarer und bestimmter nach der Ansicht des Verf. angedeutet, als er es selbst gethan hat, wodurch dann aber auch die Schwäche desselben, wie uns scheint, leichter erkannt wird. Denn offenbar ist hier ganz willkürlich angenommen, dass bei allen gleichschenkligen Dreiecken mit dem Winkel \(A\) an der Spitze und grössern Winkel an der Basis, wenn mit ihnen die im Lehrsatz angezeigte Construction vorgenommen wird, die Folge der abgeschnittenen Stücke in Rücksicht auf ihr Gleichbleiben, Grösser oder Kleinerwerden, allemal, unabhängig von der Grösse der Seiten. nothwendig dieselbe sein müsse, eine Annahme, die doch unmöglich als von selbst evident betrachtet werden darf. Da sich nun aber hierauf allein der versuchte Beweis der Unstatthaftigkeit des dritten (wie auch vierten und fünften) Falls stützt, und der ganze Artikel auch keine andere Ressourcen zum Beweise der Unstatthaftigkeit des übergangenen Falls darbietet, so glauben wir hierdurch das oben ausgesprochene Urtheil hinlänglich gerechtfertigt zu haben, wobei wir aber gern der ganzen übrigen sinnreichen Durchführung in den folgenden Artikeln volle Gerechtigkeit wiederfahren lassen.

\Needspace{10\baselineskip}
\pdfbookmark[0]{Anzeigen. Opérations géodésiques et astronomiques. 1830.}{anzeigen-operations-1830}
\begin{center}
\rule{.85\textwidth}{.3pt}\par
Göttingische gelehrte Anzeigen.\quad 1830 Februar 27.\par
\rule{.85\textwidth}{.3pt}
\end{center}
\hspace*{\parindent}\textit{Opérations géodésiques et astronomiques pour la mesure d’un arc du parallèle moyen, exécutées en Piémont et en Savoie par une commission composée d’officiers de l’état major général et d’astronomes Piémontais et Autrichiens en} 1821, 1822, 1823. \textit{Milan, de l’imprimerie impériale et royale: Tome premier} 1825. 238 S. \textit{Tome second} 1827. 412 S. in 4. Nebst einem Heft mit Figuren, Karten und sechs Rundsichten.

Die Idee der grossen Längengradmessung, von welcher die im vorliegenden Werke bekannt gemachten Operationen einen Hauptbestandtheil ausmachen, ist ursprünglich von \textsc{Laplace} ausgegangen. Seit dem Jahre 1802 waren in Oberitalien ausgedehnte Dreiecksmessungen, zunächst für topographisch-militärische
